\subsection{Varieties of sections and varieties of mappings}

15.3.1 ("Banach space structure on a space of sections"). Let B be a compact differentiable manifold of class \(C^{r}\) ; let M be a vector bundle over K (7.3.4) with base B and of class \(C^{r}\) ; denote by \(M_{r}\) the differentiable manifold of class \(C^{r}\) underlying M, and let \(\mathcal{S}_{\mathrm{M}}^{\tau}(\mathrm{B})\) be the K-vector space of sections of class \(C^{r}\) of M (7.4.1). The topology on \(\mathcal{S}_{\mathrm{M}}^{\tau}(\mathrm{B})\) induced by the topology of \(C^{r}\) -uniform convergence on \(\mathcal{C}^{\tau}(\mathrm{B};\mathrm{M}_{r})\) (12.3.10) makes \(\mathcal{S}_{\mathrm{M}}^{\tau}(\mathrm{B})\) a K-Banach space.

Let \((\mathrm{U}_{i})_{i\in\mathbb{I}}\) be a finite open covering of B and, for every \(i\in I\) , let \(c_{i}=(\mathrm{U}_{i},\varphi_{i},\mathrm{E}_{i})\) be a chart of B and \(t_{i}=(\mathrm{U}_{i},\psi_{i},\mathrm{F}_{i})\) be a K-vector chart of M (8.8.1). Let each of the \(E_{i}\) (resp. \(F_{i}\) ) be equipped with a norm compatible with its structure of Banach space over R (resp. K), and let \((\mathrm{V}_{i})_{i\in\mathbb{I}}\) be a closed covering of B, such that \(V_{i}\subset U_{i}\) for all \(i\in I\) . For \(f\in\mathcal{S}_{\mathrm{M}}^{\tau}(\mathrm{B})\) and \(i\in I\) , denote by \(f_{i}\) the map from \(\varphi_{i}(\mathrm{U}_{i})\) to \(F_{i}\) obtained by composition:

\[
\varphi_ {i} \left(\mathrm{U} _ {i}\right) \xrightarrow {\varphi_ {i} ^ {- 1}} \mathrm{U} _ {i} \xrightarrow {f} \mathrm{M} \left| \mathrm{U} _ {i} \xrightarrow {\psi_ {i}} \mathrm{U} _ {i} \times \mathrm{F} _ {i} \xrightarrow {\operatorname{pr} _ {2}} \mathrm{F} _ {i}. \right.
\]

Set

\[
\| f \| = \operatorname{Sup} _ {i \in I, p \leqslant r, x \in V _ {i}} \| D ^ {p} f _ {i} (\varphi_ {i} (x)) \|.
\]

The map \(f\mapsto\|f\|\) is a norm on \(\mathcal{S}_{\mathrm{M}}^{\tau}(\mathrm{B})\) , which gives \(\mathcal{S}_{\mathrm{M}}^{\tau}(\mathrm{B})\) the structure of a Banach space over K.

If N is an open subset of the space M, the set \(\mathcal{S}^{\tau}(\mathbf{B};\mathbf{N})\) of sections \(f\in\mathcal{S}_{\mathbf{M}}^{\tau}(\mathbf{B})\) such that \(f(\mathbf{B})\subset\mathbf{N}\) is open in \(\mathcal{S}_{\mathbf{M}}^{\tau}(\mathbf{B})\) .

15.3.2. Let B be a compact differentiable manifold of class \(C^{r}\) , and let X be a manifold of class \(C_{K}^{r,s}\) ( \(s \geqslant r + 1\) ) over B. If N is an open subset of X, denote by \(\mathcal{S}^{r}(B; N)\) the set of sections of class \(C^{r}\) of X over B whose image is contained in N. There exists on \(\mathcal{S}^{r}(B; X)\) a K-manifold structure of class \(C^{s-r}\) , and only one, such that, for every section \(\sigma \in \mathcal{S}^{r}(B; X)\) and every tubular neighborhood (M, N, \(\varphi\) ) of class \(C_{K}^{r,s}\) of \(\sigma\) (15.2.7), the triple ( \(\mathcal{S}^{r}(B; N)\) , \(\varphi^{*}\) , \(\mathcal{S}_{M}^{r}(B)\) ), where \(\varphi^{*}\) is the map \(f \mapsto \varphi \circ f\) , is a chart of the K-manifold \(\mathcal{S}^{r}(B; X)\) .

In what follows, we equip \(\mathcal{S}^{r}(\mathbf{B};\mathbf{X})\) of this structure.

15.3.3. Let X be a compact differentiable manifold of class \(C^{r}\) , and let Y be a K-manifold of class \(C^{s}(s \geqslant r + 1)\) . Let \(\mathcal{C}^{r}(X; Y)\) be the set of maps of class \(C^{r}\) from X into Y (in other words, taking values in the manifold of class \(C^{r}\) underlying Y, cf. 5.13.1 and 5.14.2). We equip \(\mathcal{C}^{r}(X; Y)\) with the structure of a K-variety of class \(C^{s-r}\) obtained by identifying it with \(\mathcal{S}^{r}(X; X \times Y)\) , cf. 15.2.3, (ii) and 15.3.2. It is called the variety of mappings of class \(C^{r}\) of X into Y. Its topology is that of \(C^{r}\) -uniform convergence (12.3.10).

Suppose Y is separated; then so is \(\mathcal{C}^{r}(X;Y)\) . Moreover, the set of immersions (resp. embeddings, submersions, surjective submersions, étale morphisms, isomorphisms) of class \(C^{r}\) from X into Y is an open subset of \(\mathcal{C}^{r}(X;Y)\) .

Let \(X'\) be a compact differentiable manifold of class \(C^r\) , \(\varphi: X' \to X\) a morphism of differentiable manifolds of class \(C^r\) and \(\psi: Y \to Y'\) a morphism of K-varieties of class \(C^s\) . The map \(f \mapsto \psi \circ f \circ \varphi\) is a morphism of class \(C^{s-r}\) of \(\mathscr{C}^r(X; Y)\) to \(\mathscr{C}^r(X'; Y')\) .

15.3.4 (“Weakening of structure”). Keep the hypotheses and notations of 15.3.3. Moreover, let \(s' \in \mathbf{N}_{\mathbf{R}}\) satisfy \(r < s' \leqslant s\) and let \(\mathbf{Y}_{s'}\) be the real manifold of class \(\mathbf{C}^{s'}\) underlying \(\mathbf{Y}\) (5.13 and 5.14.2). The structure of real manifold of class \(\mathbf{C}^{s'-r}\) on \(\mathcal{C}^r(\mathbf{X}; \mathbf{Y}_{s'})\) underlies the structure of K-manifold of class \(\mathbf{C}^{s-r}\) on \(\mathcal{C}^r(\mathbf{X}; \mathbf{Y})\) .

15.3.5. Let X and X' be compact differentiable manifolds of class \(C^{r}\) and \(C^{r'}\) , respectively, with \(r' < r\) ; let Y be a K-variety of class \(C^{s}\) , with \(s \geqslant r + 1\) . Let \(t = \operatorname{Inf}(r - r', s - r)\) . The map \((f, g) \mapsto g \circ f\) is a morphism of class \(C^{t}\) from \(\mathcal{C}^{r}(X'; X) \times \mathcal{C}^{r}(X; Y)\) to \(\mathcal{C}^{r'}(X'; Y)\) .

15.3.6 (“Tangent space to \(\mathcal{C}^r (\mathrm{X};\mathrm{Y})\)”). Given the hypotheses of 15.3.3, let \(f\in \mathcal{C}^r (\mathrm{X};\mathrm{Y})\) and let \(\xi\) be a tangent vector at \(f\) to the manifold \(\mathcal{C}^r (\mathrm{X};\mathrm{Y})\) . If \(x\in \mathrm{X}\) , denote by \(\varepsilon_{x}\) the map \(g\mapsto g(x)\) of \(\mathcal{C}^r (\mathrm{X};\mathrm{Y})\) into Y; it is a morphism of class \(\mathbf{C}^{s - r}\) . The image of \(\xi\) under \(\mathrm{T}_f(\varepsilon_x)\) is an element \(\xi_{x}\) of \(\mathrm{T}_{f(x)}(\mathrm{Y})\) ; the map \(x\mapsto \xi_{x}\) is a lifting of class \(\mathbf{C}^r\) of \(f\) in T(Y) (cf. 8.6.1) and it can be identified with a section of class \(\mathbf{C}^r\) of the vector bundle \(f^{*}\mathrm{T}(\mathrm{Y})\) of base X. In this way, one obtains an isomorphism of the tangent space \(\mathrm{T}_f(\mathcal{C}^r (\mathrm{X};\mathrm{Y}))\) onto the Banach space \(\mathcal{S}_{f^{*}\mathrm{T}(\mathrm{Y})}^r (\mathrm{X})\) .

15.3.7 ("Interpretation of \(\mathcal{C}^{p}(\mathrm{X};\mathcal{C}^{q}(\mathrm{Y};\mathrm{Z}))\)"). Let X and Y be compact differentiable manifolds of class \(\mathbf{C}^{p}\) and \(\mathbf{C}^{q}\) respectively ( \(1\leqslant p\leqslant+\infty,1\leqslant q<+\infty\) ), and let Z be a K-variety of class \(\mathbf{C}^{s}\) , with \(s\geqslant p+q\) . Let \(f\colon\mathrm{X}\times\mathrm{Y}\to\mathrm{Z}\) be a continuous map;

for every \(x \in X\) , denote by \(f_x\) the map \(y \mapsto f(x, y)\) from Y to Z. The following two properties are equivalent:

\begin{enumerate}
\def\labelenumi{\alph{enumi})}
\item
  For any charts \(c = (\mathrm{U}, \varphi, \mathrm{E})\) of \(\mathbf{X}\) , \(d = (\mathrm{V}, \psi, \mathrm{F})\) of \(\mathbf{Y}\) and \(e = (\mathrm{W}, \theta, \mathrm{G})\) of \(\mathbf{Z}\) such that \(f(\mathrm{U} \times \mathrm{V}) \subset \mathrm{W}\) , the partial derivatives \(\mathrm{D}_{\mathbb{E}}^{p'} \mathrm{D}_{\mathbb{F}}^{q'} f'\) of the expression \(f'\) of \(f\) in the charts \(c \times d\) and \(e\) (5.3.2) exist and are continuous for \(p' \leqslant p\) and \(q' \leqslant q\) .
\item
  For every \(x \in X\) , the map \(f_{x}\) is of class \(C^{q}\) and the map \(x \mapsto f_{x}\) is a morphism of class \(C^{p}\) from X into the variety \(\mathcal{C}^{q}(Y; Z)\) .
\end{enumerate}

\pandocbounded{\includegraphics[keepaspectratio]{images/part001_50717dae.jpg}}

15.3.8. Let Y be a compact K-variety of class \(C^{s}\) , with \(s \geqslant r + 1\) , and let \(\text{Diff}^{r}(Y)\) be the group of automorphisms of class \(C^{r}\) of the real manifold \(Y_{r}\) of class \(C^{r}\) underlying Y. The set \(\text{Diff}^{r}(Y)\) is open in \(\mathcal{C}^{r}(Y; Y) = \mathcal{C}^{r}(Y_{r}; Y)\) ; we endow it with the structure of a K-variety of class \(C^{s-r}\) induced by that of \(\mathcal{C}^{r}(Y; Y)\) . The composition law \((f, g) \mapsto g \circ f\) makes \(\text{Diff}^{r}(Y)\) a topological group. On the other hand, the manifold structure of class \(C^{s-r}\) is not in general compatible with the group structure of \(\text{Diff}^{r}(Y)\) . For \(f \in \text{Diff}^{r}(Y)\) , the map \(g \mapsto g \circ f\) of \(\text{Diff}^{r}(Y)\) into itself is of class \(C^{s-r}\) , but the map \(g \mapsto f \circ g\) is not in general of class \(C^{1}\) . In particular, \(\text{Diff}^{r}(Y)\) is not in general a Lie group.\footnote{Suppose that \(s = \infty\). Let \(\text{Diff}^{\infty}(Y)\) be the group of automorphisms of class \(C^{\infty}\) of Y; equip it with the coarsest topology making the canonical injections \(\text{Diff}^{\infty}(Y) \to \text{Diff}^{r}(Y)\), for every integer \(r \geqslant 1\), continuous. On the other hand, let \(\text{Diff}^{0}(Y)\) be the group of homeomorphisms of Y, endowed with the topology of uniform convergence (TG, X, § 3, no. 5, Proposition 11). If \(\dim Y \geqslant 1\), it is impossible to find a Lie group G and continuous homomorphisms from \(\text{Diff}^{\infty}(Y)\) to G and from G to \(\text{Diff}^{0}(Y)\) whose composite is the canonical injection (cf. LIE, III, § 4, Exercise 7).}

